\section{Experiencia de nivel L4: Deep Research y Claude Code}

Antes se habló de la ventana de minería; esta sección examina dónde está hoy la veta más clara. Guangmi (广密) considera que los dos mejores Agent ya ofrecen una experiencia L4: Deep Research de ChatGPT y Claude Code de Anthropic, que corresponden respectivamente a la búsqueda e investigación de información y al desarrollo de software. Hoy el mayor dividendo sigue siendo el de language/code, sobre todo el de code; todavía no es el de la multimodalidad, los modelos del mundo ni la robótica.

\begin{figure}[H]
\centering
\includegraphics[width=0.94\textwidth]{l4-experience-map.png}
\caption{Mapa de la experiencia L4: Deep Research y Claude Code representan la búsqueda y el desarrollo de software, respectivamente. Diagrama conceptual propio, elaborado a partir de la conversación de 00:44:21--00:52:43.}
\end{figure}

\begin{knowledgebox}{Lectura de la figura: por qué estos dos se parecen a L4}
Deep Research puede realizar búsquedas de información en varios pasos, sintetizarlas y redactar un informe; Claude Code puede entender un código base, modificarlo, probarlo e iterar. Todas estas tareas tienen una entrada clara, un entorno de herramientas y una salida verificable, por lo que es más fácil que den lugar a una experiencia de Agent de alto nivel.
\end{knowledgebox}

\subsection{Por qué el dividendo de code es el mayor}

El mundo del código tiene ventajas naturales: las tareas pueden formalizarse, las herramientas pueden invocarse, los resultados pueden probarse, la retroalimentación es rápida y el valor es alto. Frente a la multimodalidad, la robótica y los modelos del mundo, no sorprende que el coding haya alcanzado antes la experiencia L4. Es uno de los entornos de alto valor del mundo digital más adecuados para los Agent.

\begin{knowledgebox}{Condiciones de la experiencia L4}
\begin{tabularx}{\textwidth}{p{0.22\textwidth}p{0.34\textwidth}X}
\toprule
Condición & Deep Research & Claude Code \\
\midrule
Tarea clara & Preguntas de investigación y recopilación de fuentes & issue, bug, feature. \\
Herramientas disponibles & Búsqueda, páginas web, citas & Código base, pruebas, shell. \\
Resultado evaluable & Calidad del informe y de las fuentes & Pruebas aprobadas, review del código. \\
Valor alto & Ahorra tiempo de investigación & Ahorra tiempo de ingeniería. \\
\bottomrule
\end{tabularx}
\end{knowledgebox}

\begin{knowledgebox}{Próximos escenarios candidatos a L4}
\begin{tabularx}{\textwidth}{p{0.20\textwidth}p{0.34\textwidth}X}
\toprule
Escenario & Por qué podría surgir L4 & Dificultad \\
\midrule
Legal/cumplimiento & Uso intensivo de documentos, procesos claros, valor alto & Gran responsabilidad, alto costo de los errores. \\
Investigación financiera & Recuperación de información, modelos, informes y pistas de operaciones & Altas exigencias de licencias de datos y de tiempo real. \\
Operaciones de ventas & CRM, correo, reuniones y cotizaciones pueden convertirse en herramientas & Datos privados y procesos organizacionales complejos. \\
Análisis de datos & Ciclo cerrado de SQL, reportes, BI e interpretación del negocio & Requiere entender la definición de las métricas y la semántica del negocio. \\
Diseño/multimodal & La generación visual y la iteración de edición se aceleran & La evaluación subjetiva y los derechos de autor son más complejos. \\
\bottomrule
\end{tabularx}
\end{knowledgebox}

\subsection{Resumen de la sección}

La experiencia L4 no es un nivel abstracto, sino que el usuario esté realmente dispuesto a confiarle tareas complejas al sistema. Deep Research y Claude Code son hoy las dos muestras más claras.
